\documentclass[11pt]{article}

\usepackage[margin=1in]{geometry}
\usepackage{amsmath,amssymb,amsthm}
\usepackage{graphicx}
\usepackage[hidelinks]{hyperref}

\newtheorem{theorem}{Theorem}
\newtheorem{remark}{Remark}

\newcommand{\Bess}{\mathrm{Bess}}
\newcommand{\eps}{\varepsilon}
\newcommand{\norm}[1]{\left\lVert #1 \right\rVert}

\title{A Non-Uniform Crossnorm Partial Answer to an Unconditional-Basis Tensor Product Problem}
\author{Automated Math Discovery Run \texttt{fa\_banach\_001}}
\date{June 12, 2026}

\begin{document}
\maketitle

\section*{Source and Status}

This packet concerns the tensor-norm question in:

\begin{quote}
Rafik Karkri and Samir Kabbaj,
\emph{A New Crossnorm That Preserves Unconditional Bases in Banach Spaces},
arXiv:2506.22522.
\end{quote}

The source paper frames the original problem as the existence of a tensor
norm, that is, a uniform reasonable crossnorm, preserving unconditional
Schauder bases for all Banach-space tensor products. Page 1 also states
that this uniform existence question remains open, while the paper
constructs a generally non-uniform crossnorm for each fixed pair.

\begin{center}
\includegraphics[width=\textwidth]{figures/open_problem_crop.png}
\end{center}

\section*{Literature-Implied Partial Result}

\begin{theorem}
Let \(E\) and \(F\) be Banach spaces equipped with unconditional Schauder
bases
\[
  \mathcal F=((a_m,b_m^*))_{m\ge 1},\qquad
  \mathcal G=((x_n,y_n^*))_{n\ge 1}.
\]
For this fixed pair of spaces and chosen bases, the source paper
constructs equivalent Besselian norms on \(E\) and \(F\) and a
pair-and-basis-dependent crossnorm
\[
  \alpha_{\mathcal F,\mathcal G}^{\Bess}
\]
on \(E\otimes F\) such that:
\begin{enumerate}
\item \(\alpha_{\mathcal F,\mathcal G}^{\Bess}\) is a reasonable
crossnorm on the tensor product of the corresponding Besselian
renormings, and
\[
  \eps_{\mathcal F,\mathcal G}(u)
  \le
  \alpha_{\mathcal F,\mathcal G}^{\Bess}(u)
  \le
  \pi_{\mathcal F,\mathcal G}(u)
  \qquad (u\in E\otimes F).
\]
\item The square-ordered tensor product sequence
\[
  \mathcal F\otimes\mathcal G
  =
  ((a_m\otimes x_n,b_m^*\otimes y_n^*))_{m,n\ge 1}
\]
forms an unconditional Schauder basis of
\[
  E\otimes_{\alpha_{\mathcal F,\mathcal G}^{\Bess}} F.
\]
\end{enumerate}
\end{theorem}

\begin{proof}
This is a source-backed literature-status statement rather than an
independent new proof.

The construction of
\(\alpha_{\mathcal F,\mathcal G}^{\Bess}\) is given in the source paper by
the definition preceding Proposition \(\texttt{New-crossnorm}\). That
proposition proves that, after passing to the equivalent Besselian norms
on \(E\) and \(F\), the constructed functional is a reasonable crossnorm
and satisfies the injective/projective bounds displayed above.

When \(\mathcal F\) and \(\mathcal G\) are unconditional Schauder bases,
the corresponding source theorem applies directly; in the source file it
is labelled \(\texttt{tensor.bess.is.bess}\). It states that the tensor
product sequence, ordered by square ordering, forms an unconditional
Schauder basis in the completion for
\(\alpha_{\mathcal F,\mathcal G}^{\Bess}\). This is exactly the asserted
non-uniform preservation statement.
\end{proof}

\section*{Why This Is Only Partial}

The original problem asks for a tensor norm in the usual sense: a uniform
reasonable crossnorm defined functorially across Banach-space tensor
products. The Besselian crossnorm above depends on the particular pair
\((E,F)\) and on the chosen bases \((\mathcal F,\mathcal G)\). The source
paper also states and demonstrates that this crossnorm is generally not
uniform.

Therefore the packet records a strong partial answer:
\[
  \text{pair/basis-dependent reasonable crossnorm preserving USBs,}
\]
not a full solution of the uniform tensor-norm existence problem.

\section*{Review Checklist}

\begin{itemize}
\item Check that the downstream problem accepts equivalent Besselian
renormings; otherwise the result should be cited with this renorming
qualification.
\item Check that the term ``tensor norm'' is not used for this
construction without the non-uniform caveat.
\item Treat this packet as a literature-status partial result unless a
separate argument upgrades the construction to a uniform tensor norm.
\end{itemize}

\end{document}
